\section{A finite Euler-sum space for every centered harmonic moment}
\label{hpnew:section}

The recent raw-power report proves that the centered series
\[
 F_p(s,a)=\sum_{n=0}^{\infty}\frac{H_n^p}{(n+a)^s},\qquad
 H_n=\sum_{j=1}^n\frac1j,\quad H_0=0,
\]
is regular at $s=0,-2,-4,\ldots$ when $a=1/2$.
Its Question 10 asks for a generator that avoids listing every
composition of $p$, with the complete sequence $F_p(-2,1/2)$ as a
first target \cite{RepoExactResonant}. We give such a generator for
\emph{every} nonnegative even moment. More strongly, for fixed $p$ all
these values lie in one finite rational span of convergent Euler sums,
independent of the moment order. These are exact analytic identities;
no assertion of arithmetic independence or literature-wide priority is
made. They concern values at regular spectral points, not derivatives
of a formula known only at those points.

\subsection{An exponent parameter with the required joint regularity}

Introduce the auxiliary exponent family
\begin{equation}\label{hpnew:exponent}
 \mathscr F(s,u)=\sum_{n=0}^{\infty}
       \frac{\exp(uH_n)}{(n+1/2)^s},\qquad
                    \Re(s-u)>1.
\end{equation}
Positive bases have their real logarithms. This normally convergent
series is jointly holomorphic in its initial domain. At $u=0$ its
$p$th derivative is $F_p(s,1/2)$ for $p\ge1$; the zeroth derivative is
$\zeta(s,1/2)$. Let $B_j(x)$ denote the Bernoulli polynomials, and put
\begin{equation}\label{hpnew:coefficients}
 b_j=-\frac{B_{2j}(1/2)}{2j},\qquad
 c_j(u)=[z^j]\exp\left(u\sum_{\ell\ge1}b_\ell z^\ell\right).
\end{equation}
The latter definition is formal: each $c_j$ is a finite polynomial.
No convergence of the infinite Bernoulli series is required.

\begin{lemma}[Joint continuation and the centered divisor pattern]
\label{hpnew:joint}
The family $\mathscr F$ has a joint meromorphic continuation to
$\mathbb C^2$, with possible polar hyperplanes
\begin{equation}\label{hpnew:divisors}
                         s-u=1-2j,\qquad j=0,1,2,\ldots.
\end{equation}
In particular it is jointly holomorphic near every $(-2r,0)$,
$r\ge0$. The function
\begin{equation}\label{hpnew:V}
 V_r(u)=\mathscr F(-2r,u)
       =\sum_{p\ge1}F_p(-2r,1/2)\frac{u^p}{p!}
\end{equation}
is holomorphic for $|u|<1$.
\end{lemma}

\begin{proof}
With $x=n+1/2$, Euler--Maclaurin gives, for every fixed $J$,
\[
 H_n=\log x+\gamma+\sum_{j=1}^Jb_jx^{-2j}
                              +O(x^{-2J-2}).
\]
This is the centered shifted-digamma expansion used in
\cite{RepoExactResonant}. Exponentiating a finite truncation gives
\begin{equation}\label{hpnew:asymptotic}
 \exp(uH_n)=e^{\gamma u}x^u
      \left(\sum_{j=0}^Jc_j(u)x^{-2j}
                         +O_K(x^{-2J-2})\right),
\end{equation}
uniformly for $u$ in a compact set $K$. Fixed $u$ derivatives have
analogous estimates with powers of $\log x$.
Subtract these finitely many terms in \eqref{hpnew:exponent}, starting
at any fixed positive index. The remainder is normally convergent
when $\Re(s-u)>-2J-1$, and the subtracted terms are
$e^{\gamma u}c_j(u)\zeta(s-u+2j,N+1/2)$.
The sole pole of each Hurwitz function is at its first argument $1$
\cite{DLMFHurwitz}. Arbitrarily large $J$ prove the stated joint
continuation. None of the hyperplanes \eqref{hpnew:divisors} passes
through $(-2r,0)$. Joint holomorphy, rather than a formal interchange
at a singular point, therefore identifies all the coefficients in
\eqref{hpnew:V}. Its constant coefficient is
$\zeta(-2r,1/2)=0$.
\end{proof}

\subsection{A triangular generator from ordinary telescoping}

For $k\ge2$ define the normally convergent functions
\begin{align}
 E_k(u)&=\sum_{n\ge1}\frac{\exp(uH_{n-1})}{n^k},
                                  &&\Re u<k-1,\label{hpnew:Ek}\\
 B(u)&=1+\sum_{k\ge2}\frac{u^k}{k!}E_k(u),
                                  &&\Re u<1.\label{hpnew:B}
\end{align}
The letter $B$ without an index in \eqref{hpnew:B} is an auxiliary
function, not a Bernoulli polynomial. Its coefficients involve the
ordinary convergent Euler sums
\begin{equation}\label{hpnew:Euler}
 E_{q,k}=\sum_{n\ge1}\frac{H_{n-1}^q}{n^k},\qquad
 q\ge0,\quad k\ge2;\qquad
 E_k(u)=\sum_{q\ge0}E_{q,k}\frac{u^q}{q!}.
\end{equation}
The Taylor identity in \eqref{hpnew:Euler} is used near $u=0$.
Here $H_{n-1}^0=1$, including at $n=1$.

\begin{theorem}[A finite triangular generator at every centered moment]
\label{hpnew:generator}
For $m=0,1,\ldots$, define $G_m(u)$ recursively by
\begin{align}
 (u+m+1)G_m(u)
 ={}&\sum_{j=0}^{m-1}
 \left\{(-1)^{m+1-j}\binom{m+1}{j}
             -\frac{u^{m+1-j}}{(m+1-j)!}\right\}G_j(u)\notag\\
 &+\frac{u^{m+1}}{(m+2)!}B(u)
 -\sum_{k\ge2}\frac{u^{k+m+1}}{(k+m+1)!}E_k(u).
 \label{hpnew:Grecurrence}
\end{align}
The first sum is empty at $m=0$. All functions are holomorphic on
$|u|<1$, and
\begin{equation}\label{hpnew:Vrecurrence}
 \boxed{\displaystyle
 V_r(u)=\sum_{j=0}^{2r}\binom{2r}{j}
                     \left(-\frac12\right)^{2r-j}G_j(u).}
\end{equation}
Thus $2r+1$ triangular steps determine the exponential generating
function for all $F_p(-2r,1/2)$. For a specified $p$, only finitely
many coefficients of finitely many $E_k$ are needed.
\end{theorem}

\begin{proof}
Put $g_n(u)=\exp(uH_{n-1})$. It obeys the exact discrete equation
\[
                  g_{n+1}(u)=e^{u/n}g_n(u).
\]
For $\Re u<0$, telescoping and then an absolutely convergent
exponential expansion give
\[
 -1=\sum_{n\ge1}(g_{n+1}-g_n)
   =u\sum_{n\ge1}\frac{g_n}{n}
                      +\sum_{k\ge2}\frac{u^k}{k!}E_k(u).
\]
Consequently the continued order-one sum is exactly
\begin{equation}\label{hpnew:Eone}
                \sum_{n\ge1}\frac{g_n(u)}n=-\frac{B(u)}u.
\end{equation}
This equality initially concerns ordinary sums in $\Re u<0$.

For a fixed $m\ge0$, work first in $\Re u<-m-1$, where all the
following sums converge absolutely. Multiplying the discrete equation
by $n^{m+1}$, summing to a finite cutoff $N$, and only then
shifting an index yields the following formulas after $N\to\infty$
\begin{align*}
 \sum_{n\ge1}n^{m+1}(g_{n+1}-g_n)
 &=\sum_{j=0}^{m}(-1)^{m+1-j}\binom{m+1}{j}
                                  \sum_{n\ge1}n^jg_n,\\
 \sum_{n\ge1}n^{m+1}(g_{n+1}-g_n)
 &=\sum_{k\ge1}\frac{u^k}{k!}
                                  \sum_{n\ge1}\frac{g_n}{n^{k-m-1}}.
\end{align*}
The upper-end term is $N^{m+1}g_{N+1}=O(N^{m+1+\Re u})$ and
vanishes in the stated half-plane. There is no lower-end boundary
term because the polynomial being shifted has positive degree.
This finite-cutoff argument never separates a difference into two
possibly divergent infinite sums. In the second line separate $k=1$,
$2\le k\le m+1$, $k=m+2$, and $k\ge m+3$.
Using \eqref{hpnew:Eone} for the order-one term and moving the two
copies of $\sum n^mg_n$ to the left gives
\eqref{hpnew:Grecurrence}. In particular its $G_m$ is the
meromorphic continuation of the ordinary sum
$\sum_{n\ge1}n^m\exp(uH_{n-1})$ from $\Re u<-m-1$.

Here are the convergence details for the continuation. On a compact
subset of $\Re u<1$, choose $U<1$ with $\Re u\le U$.
The two-sided estimate $H_{n-1}=\log n+O(1)$ gives
$|g_n(u)|\le Cn^U$, uniformly, including when $\Re u$ is negative. Hence all $E_k(u)$, $k\ge2$, are
bounded by a common summable majorant, and the factorial denominators
make both infinite $k$ sums normally convergent. Fixed derivatives in
$u$ are justified on slightly larger compact subsets, or directly by
inserting powers of $H_{n-1}$. The recurrence therefore continues
$G_m$ meromorphically to $\Re u<1$, with possible poles only at
$-1,-2,\ldots,-m-1$. In particular it is holomorphic on $|u|<1$.

Finally, for $\Re u<-2r-1$, an ordinary change of index in the
absolutely convergent defining series gives
\[
 \mathscr F(-2r,u)=\sum_{n\ge1}(n-1/2)^{2r}g_n(u)
   =\sum_{j=0}^{2r}\binom{2r}{j}(-1/2)^{2r-j}G_j(u).
\]
Meromorphic continuation in $u$ proves \eqref{hpnew:Vrecurrence};
Lemma~\ref{hpnew:joint} identifies its Taylor coefficients with the
claimed regular spectral values.
\end{proof}

\begin{remark}[Why the auxiliary values cannot be specialized in the other order]
\label{hpnew:orderwarning}
Induction in \eqref{hpnew:Grecurrence} gives $G_m(0)=0$ for every
$m$. For example $G_1(0)=0$, whereas $\zeta(-1)=-1/12$.
This is consistent: $G_m$ first fixes the spectral argument at $-m$
and then continues the exponent $u$. The uncentered family has polar hyperplanes $s-u=1-j$.
Even a single such divisor can prevent joint holomorphy when its
numerator vanishes at the specialization point. For example its
local subtraction includes
\[
 e^{\gamma u}\left(\frac{u^2}{8}-\frac{u}{12}\right)
                                      \zeta(s-u+2).
\]
At $s=-1$ and then $u\to0$, this contributes $1/12$; at $u=0$
first, it contributes zero. It is therefore invalid to replace
$G_m(0)$ by the separately specialized zeta value.
The centered combination \eqref{hpnew:Vrecurrence} is justified by
joint holomorphy, which was proved before coefficient extraction.
\end{remark}

\subsection{A finite arithmetic space independent of the moment order}

\begin{theorem}[Uniform finite Euler-sum span]
\label{hpnew:uniformspan}
For $p\ge1$, set
\begin{equation}\label{hpnew:space}
 \mathcal E_{p-1}=\operatorname{span}_{\mathbb Q}
 \left(\{1\}\cup
   \{E_{q,k}:q\ge0,\ k\ge2,\ q+k\le p-1\}\right).
\end{equation}
Then
\begin{equation}\label{hpnew:spanconclusion}
 \boxed{\displaystyle
             F_p(-2r,1/2)\in\mathcal E_{p-1}
                    \quad\hbox{for every }r\ge0.}
\end{equation}
Every value has a finite multiple-zeta expression using weights at
most $p-1$, with no Euler constant required. In particular every
centered even-negative moment of $H_n$ or $H_n^2$ is rational.
For $p\ge2$, the displayed spanning list has at most
$1+\binom{p-1}{2}$ elements; this is a spanning bound, not a claim
that the periods in the list are independent.
\end{theorem}

\begin{proof}
The recurrence is linear in the functions $E_k(u)$, and all rational
coefficients being expanded are regular at $u=0$.
An occurrence of $E_k$ first has a factor $u^{k+m+1}$, either in
the last sum of \eqref{hpnew:Grecurrence} or through $B$ in its
preceding term. Its exponent is at least $k+1$.
Induction on $m$ therefore shows that the coefficient of $u^p$ in
any $G_m$, and hence in any $V_r$, is a rational linear combination
of $1$ and $E_{q,k}$ with $q+k\le p-1$.

For completeness, the finite multiple-zeta reduction is
\begin{equation}\label{hpnew:MZV}
 E_{q,k}=\sum_{\substack{(r_1,\ldots,r_d)\text{ a composition of }q}}
       \frac{q!}{r_1!\cdots r_d!}\,
                          \zeta(k,r_1,\ldots,r_d).
\end{equation}
At $q=0$ the empty composition gives $E_{0,k}=\zeta(k)$.
To prove the formula, expand $H_{n-1}^q$, order its distinct indices,
and group equal indices into blocks of sizes $r_1,\ldots,r_d$.
The displayed multinomial is their multiplicity. Each resulting MZV
converges because its first index is $k\ge2$, and its weight is
$q+k$. Formula \eqref{hpnew:MZV} is an optional final conversion:
it is not necessary for computing the generator itself.
\end{proof}

This uniform space is smaller information than an unspecified
multiple-zeta algebra at each separate moment. Increasing $r$ changes
rational coefficients in one fixed finite list; it does not require
Euler sums of increasing weight.

\subsection{Two all-moment Bernoulli identities}

\begin{corollary}[The first two powers at every centered even moment]
\label{hpnew:Bernoulli}
For every $r\ge0$,
\begin{align}
 F_1(-2r,1/2)&=\frac12 B_{2r}(1/2),\label{hpnew:firstpower}\\
 F_2(-2r,1/2)&=-\sum_{j=0}^{r}\binom{2r}{2j}
       \frac{B_{2j}B_{2r-2j}(1/2)}{2j+1}.
                         \label{hpnew:secondpower}
\end{align}
\end{corollary}

\begin{proof}
The classical generating functions needed here are particularly short.
Write $L(t)=-\log(1-e^{-t})$, and define
$A_p(t)=(1-e^{-t})\sum_{n\ge0}H_n^pe^{-nt}$.
Then $A_1=L$ and $A_2=L^2+\Li_2(e^{-t})$.
Consequently
\[
                         A_2'(t)=-\coth(t/2)L(t).
\]
The odd powers in the local expansion of $L$ consist only of $t/2$:
\[
 L(t)=-\log t+\frac t2
                  -\log\left(\frac{\sinh(t/2)}{t/2}\right).
\]
Here ``odd powers'' refers to the power--log expansion at $t=0$,
with no branch continuation to negative $t$ being assumed.
It follows that the odd-power part of $A_2$ is
\[
 -\int_0^t\frac v2\coth(v/2)\,dv
       =-\sum_{j\ge0}\frac{B_{2j}}{(2j)!(2j+1)}t^{2j+1}.
\]
The centered Mellin kernel is $A_p(t)/(2\sinh(t/2))$.
Its log-free coefficient of $t^{2r}$, multiplied by $(2r)!$, is
$F_p(-2r,1/2)$. Indeed that coefficient gives the simple Mellin
pole at $s=-2r$, and the simple zero of $1/\Gamma(s)$ cancels
it; all holomorphic Mellin terms vanish at the specialization.
Regularity proved in Lemma~\ref{hpnew:joint} eliminates higher
logarithmic poles. Finally use
\[
 \frac1{2\sinh(t/2)}=
             \sum_{j\ge0}\frac{B_{2j}(1/2)}{(2j)!}t^{2j-1}.
\]
Multiplying the two convergent local series gives
\eqref{hpnew:firstpower} and \eqref{hpnew:secondpower}.
\end{proof}

\subsection{A sufficient parity class for generalized harmonic products}

There is also a direct extension toward Question 11 of
\cite{RepoExactResonant}. Define
$H_n^{(q)}=\sum_{k=1}^n k^{-q}$ and
$U_q(s)=\sum_{n\ge0}H_n^{(q)}(n+1/2)^{-s}$.

\begin{corollary}[Odd harmonic orders and finite ordinary-zeta values]
\label{hpnew:oddorders}
Let $P$ be any polynomial in finitely many variables. The Dirichlet
series with numerator
\[
                 P(H_n,H_n^{(3)},H_n^{(5)},\ldots,H_n^{(2d+1)})
\]
is holomorphic at every $s=-2r$, $r\ge0$. This is a sufficient
class of numerator polynomials, not a necessity classification.
For a single positive odd integer $q$, the exact values are
\begin{equation}\label{hpnew:oddformula}
 \boxed{\displaystyle
 U_q(-2r)=-\frac1{2r+1}\sum_{j=0}^{r}\binom{2r+1}{2j}
          B_{2j}(1/2)\,
                         \zeta(q-2r-1+2j).}
\end{equation}
For fixed odd $q\ge3$, all these moments belong to the same space
$\operatorname{span}_{\mathbb Q}\{1,\zeta(2),\zeta(4),\ldots,
\zeta(q-1)\}$. At $q=1$, formula \eqref{hpnew:oddformula} reduces
to \eqref{hpnew:firstpower}.
\end{corollary}

\begin{proof}
For $q\ge3$ and $x=n+1/2$, use
$H_n^{(q)}=\zeta(q)-\zeta(q,x+1/2)$ and the midpoint expansion
\[
 \zeta(q,x+1/2)=\frac{x^{1-q}}{q-1}
 +\sum_{j=1}^{J}\frac{B_{2j}(1/2)(q)_{2j-1}}{(2j)!}
                                       x^{1-q-2j}
 +O(x^{-q-2J-1}).
\]
When $q$ is odd, every displayed inverse power is even. Together
with the even inverse-power expansion of $H_n$ following its
logarithm, this shows that any stated polynomial has an asymptotic
expansion containing only $x^{-2j}(\log x)^h$. Finite subtraction
in its Dirichlet series leaves possible poles only at $s=1-2j$.
This proves regularity at every nonpositive even integer.

To prove the value formula, first take $\Re q>2r+2$. The identity
$U_q(s)-\zeta(q)\zeta(s,1/2)
=-\sum_{n\ge0}\zeta(q,n+1)(n+1/2)^{-s}$ continues the difference
to $s=-2r$ by an absolutely convergent tail sum. Since
$\zeta(-2r,1/2)=0$, ordinary summation reversal gives
\[
 U_q(-2r)=-\sum_{k\ge1}k^{-q}
                         \sum_{n=0}^{k-1}(n+1/2)^{2r}.
\]
The centered Faulhaber identity is
\[
 \sum_{n=0}^{k-1}(n+1/2)^{2r}
 =\frac{B_{2r+1}(k+1/2)}{2r+1}
 =\frac1{2r+1}\sum_{j=0}^{r}\binom{2r+1}{2j}
                              B_{2j}(1/2)k^{2r+1-2j}.
\]
Insertion proves \eqref{hpnew:oddformula} in this half-plane.

The same tail expansion supplies joint meromorphic continuation
in $(s,q)$; its possible moving divisors are $s+q=2-2j$.
None passes through $(-2r,q)$ for an odd positive integer $q$.
The apparent pole at $q=1$ cancels: near $s=-2r$ the combination
\[
 \zeta(q)\zeta(s,1/2)-\frac{\zeta(s+q-1,1/2)}{q-1}
\]
is jointly holomorphic after writing it as
$(\zeta(q)-(q-1)^{-1})\zeta(s,1/2)
-(\zeta(s+q-1,1/2)-\zeta(s,1/2))/(q-1)$.
The remaining finite subtractions and the normal remainder have
no divisor there. Thus meromorphic continuation of the value formula
really gives the separately specified odd-order functions, including
$H_n^{(1)}=H_n$. Finally, every zeta argument in
\eqref{hpnew:oddformula} is even; negative even arguments give zero
and $\zeta(0)=-1/2$, proving the claimed fixed span.
\end{proof}

For example,
\begin{align}\label{hpnew:thirdorder}
 U_3(0)&=-\zeta(2),&
 U_3(-2)&=\frac16+\frac1{12}\zeta(2),&
 U_3(-4)&=-\frac1{12}-\frac7{240}\zeta(2),\\
 U_5(0)&=-\zeta(4),&
 U_5(-2)&=-\frac13\zeta(2)+\frac1{12}\zeta(4),\notag\\
 U_5(-4)&=\frac1{10}+\frac16\zeta(2)-\frac7{240}\zeta(4).
                                                               \notag
\end{align}
The odd-order restriction should be retained. At an even-order
specialization a vanishing coefficient can still meet a moving divisor,
so continuing the fixed spectral value in $q$ need not agree with
specializing $q$ before taking the spectral value. No such ambiguity
occurs at the odd orders proved here.

\subsection{Compact generators at zero, minus two, and minus four}

Put $Q(u)=V_0(u)$, $V(u)=V_1(u)$, and $W(u)=V_2(u)$.
The first three moment rows have the following compact forms.
\begin{corollary}[Three complete rows in the harmonic power]
\label{hpnew:threerows}
For $|u|<1$,
\begin{align}
 (u+1)Q(u)
 &=\frac u2+\frac12\sum_{k\ge2}(k-1)
                  \frac{u^{k+1}}{(k+1)!}E_k(u),\label{hpnew:Q}\\
 (u+3)V(u)
 &=\frac{u^3-3u-3}{12}Q(u)-\frac{u^3}{24}\notag\\
 &\hspace{4mm}-\frac1{24}\sum_{k\ge2}(k-1)(k+1)(k+6)
                  \frac{u^{k+3}}{(k+3)!}E_k(u),\label{hpnew:minus2}\\
 (u+5)W(u)
 &=-\frac{u^3+18u^2+90u+150}{60}V(u)\notag\\
 &\hspace{4mm}-\frac{u^3+6u^2+15u+15}{240}Q(u)\notag\\
 &\hspace{4mm}+\frac1{120}\sum_{k\ge2}k(k-1)(k+1)
                  \frac{u^{k+5}}{(k+5)!}E_k(u).
                                      \label{hpnew:minus4}
\end{align}
\end{corollary}

\begin{proof}
The $m=0$ case of \eqref{hpnew:Grecurrence}, followed by substitution
of \eqref{hpnew:B}, is \eqref{hpnew:Q}. For the next row use
\[
 V=G_2-G_1+\frac14G_0.
\]
The $G_1$ coefficient is $-u(u+2)/2$ before its substitution, so its
apparent denominator $u+2$ cancels. Substituting $B$ gives the factor
\[
 \frac{k+1}{2(k+3)!}-\frac1{24k!}
       =-\frac{(k-1)(k+1)(k+6)}{24(k+3)!},
\]
which proves \eqref{hpnew:minus2}.
For the last row substitute $m\le4$ into
\[
 W=G_4-2G_3+\frac32G_2-\frac12G_1+\frac1{16}G_0.
\]
Collect $Q$, $V$, the constant part of $B$, and the coefficient of
each $E_k$ separately. The first three are respectively
$-(u^3+6u^2+15u+15)/(240(u+5))$,
$-(u^3+18u^2+90u+150)/(60(u+5))$, and zero; the last is
$k(k-1)(k+1)u^{k+5}/(120(u+5)(k+5)!)$.
These rational simplifications prove \eqref{hpnew:minus4}. The
verification script replays them with symbolic $u$ and $k$.
\end{proof}

The zero row \eqref{hpnew:Q} recovers the recurrence already proved
in \cite{RepoExactResonant}. Equations \eqref{hpnew:minus2} and
\eqref{hpnew:minus4} determine two further complete rows, while
Theorem~\ref{hpnew:generator} supplies every remaining row.
For an explicit finite coefficient version, write
$Q_p=F_p(0,1/2)$, $V_p=F_p(-2,1/2)$ and $Q_0=V_0=0$.
For $p\ge1$, \eqref{hpnew:minus2} says
\begin{align}\label{hpnew:finiteV}
 3V_p+pV_{p-1}
 ={}&\frac{p(p-1)(p-2)}{12}Q_{p-3}
       -\frac p4Q_{p-1}-\frac14Q_p-\frac14\mathbf1_{p=3}\notag\\
 &-\frac1{24}\sum_{k=2}^{p-3}(k-1)(k+1)(k+6)
                     \binom p{k+3}E_{p-k-3,k}.
\end{align}
The first term is omitted for $p<3$, and an empty sum is zero.
There is no infinite coefficient calculation at a fixed $p$.

For example, the next fifth-power identity is
\begin{equation}\label{hpnew:fifth}
 \boxed{\displaystyle
 F_5(-2,1/2)=-\frac{200}{81}-\frac{23}{54}\zeta(2)
                  +\frac5{12}\zeta(3)-\frac{11}{8}\zeta(4).}
\end{equation}
At the next spectral point the first five powers are
\begin{align}
 F_1(-4,1/2)&=\frac7{480},&
 F_2(-4,1/2)&=\frac{19}{3600},\notag\\
 F_3(-4,1/2)&=\frac{493}{18000}+\frac7{480}\zeta(2),
                                                      \label{hpnew:table4}\\
 F_4(-4,1/2)&=-\frac{6479}{67500}+\frac{19}{1800}\zeta(2)
                              +\frac7{80}\zeta(3),\notag\\
 F_5(-4,1/2)&=\frac{98437}{202500}+\frac{643}{5400}\zeta(2)
                +\frac{19}{240}\zeta(3)+\frac{77}{160}\zeta(4).
                                                        \notag
\end{align}
Coefficient extraction from \eqref{hpnew:Q}--\eqref{hpnew:minus4}
proves these evaluations. Beyond $E_{0,k}=\zeta(k)$, the only
low-weight reductions needed are the classical identities
$E_{1,2}=\zeta(3)$, $E_{1,3}=\zeta(4)/4$ and
$E_{2,2}=11\zeta(4)/4$, also used in \cite{RepoExactResonant}.

\subsection{Exact exponent singularities and reproducible checks}

The apparent negative even denominators in the triangular recursion
cancel in every centered row. This has a stronger analytic explanation.
\begin{corollary}[Exact negative poles of the exponent generator]
\label{hpnew:upoles}
On $\Re u<1$, the function $V_r$ has exactly the simple poles
\[
                     u=-2r-1,-2r+1,\ldots,-1.
\]
At $u_j=2j-2r-1$, $0\le j\le r$, its residue is
\begin{equation}\label{hpnew:residue}
             \operatorname{Res}_{u=u_j}V_r(u)
                        =-e^{\gamma u_j}c_j(u_j).
\end{equation}
Its Taylor series at zero therefore has exact radius $1$.
\end{corollary}

\begin{proof}
The finite subtraction in \eqref{hpnew:asymptotic} leaves precisely
these possible poles in the indicated half-plane. The Hurwitz pole is
$\zeta(1-(u-u_j),N+1/2)$, with residue $-1$ in $u$, proving
\eqref{hpnew:residue}. To check that none is removable, use
$B_{2j}(1/2)=(2^{1-2j}-1)B_{2j}$. Thus
$b_j=(-1)^{j+1}|b_j|$. For a negative real $u$, every term in the
coefficient $c_j(u)$ has sign $(-1)^j$, and at least one term is
nonzero; also $c_0=1$. Hence every displayed residue is nonzero.
The pole at $-1$ fixes the exact Taylor radius.
\end{proof}

The accompanying \src{code/verify_harmonic_powers.py} performs exact
rational checks of the three generator rows, extracts the displayed
values, checks the first two Bernoulli formulas and the centered
Faulhaber identity through multiple moment orders, and evaluates six
odd generalized-harmonic checkpoints. Its independent numerical
calculation starts from the
original harmonic-power sums, subtracts their centered Euler--Maclaurin
expansion, and evaluates the remaining Hurwitz order derivatives.
Two different truncation lengths are compared. This diagnostic tests
signs and constants without reusing the triangular generator; it is
not a proof of an infinite identity or a certified numerical error bound.
The proof is the normally convergent telescoping argument and the
joint continuation established above. The first regular spectral
coefficient beyond these centered values still carries global Mellin
information; the theorem makes no unsupported claim to eliminate it.
